\documentclass[10pt,a4paper]{article}
\usepackage[utf8]{inputenc}
\usepackage[spanish,es-tabla]{babel}
\usepackage{geometry}
\geometry{top=2.0cm, bottom=2.0cm, left=2.0cm, right=2.0cm}
\usepackage{amsmath,amssymb}
\usepackage{booktabs}
\usepackage{array}
\usepackage{tabularx}
\usepackage{xcolor}
\usepackage{enumitem}
\usepackage{fancyhdr}
\usepackage{titlesec}
\usepackage{tcolorbox}
\usepackage{pifont}

\definecolor{tecblue}{RGB}{0, 43, 91}
\definecolor{alertred}{RGB}{180, 0, 0}
\definecolor{okgreen}{RGB}{0, 120, 40}

\pagestyle{fancy}
\fancyhf{}
\lhead{\textbf{EL-2113 Circuitos en CC} -- Protocolo de Pruebas de Laboratorio}
\rhead{Sistema de Comunicación Óptica}
\cfoot{\thepage}

\titleformat{\section}{\large\bfseries\color{tecblue}}{\thesection.}{0.5em}{}[\titlerule]
\titleformat{\subsection}{\normalsize\bfseries\color{tecblue}}{\thesubsection.}{0.5em}{}

\newcommand{\cmark}{\textcolor{okgreen}{\ding{51}}}
\newcommand{\xmark}{\textcolor{alertred}{\ding{55}}}
\newcommand{\squarebox}{\raisebox{-0.1ex}{\Large $\square$}}

\begin{document}

\begin{center}
    {\LARGE \textbf{\color{tecblue}PROTOCOLO MAESTRO DE PRUEBAS DE LABORATORIO}}\\[0.2cm]
    {\large \textbf{Validación Experimental y Cumplimiento Normativo (Req. A.1--Q.18)}}\\[0.15cm]
    \textbf{Curso:} Circuitos Eléctricos en Corriente Continua (EL-2113) \quad|\quad \textbf{Fecha:} Octubre 2026\\[0.1cm]
    \textbf{Escuela de Ingeniería Electrónica -- Instituto Tecnológico de Costa Rica (ITCR)}\\[0.1cm]
    \textbf{Equipo de Trabajo:} Mario Cerdas (2022437519, m.cerdas@estudiantec.cr), Sebasti\'an Quir\'os Fern\'andez (2024073526, s.quiros.4@estudiantec.cr), Ian Fuentes N\'u\~nez (2025100912, I.fuentes.1@estudiantec.cr)\\[0.1cm]
    \textbf{Base F\'isica:} Maqueta Corte L\'aser LAIMI + Kit 3D M3 \quad|\quad \textbf{Estado:} Ensamble 100\% completado, puentes LM358N corregidos y validado
\end{center}

\vspace{-0.2cm}
\begin{tcolorbox}[colback=tecblue!5, colframe=tecblue, title=\textbf{OBJETIVO DE LA SESIÓN EXPERIMENTAL}]
Validar de forma secuencial y cuantitativa las 7 etapas del prototipo físico del sistema de comunicación óptica, verificar el margen de error menor al $10\%$ en el convertidor DAC (Req. O.16), comprobar la modulación acústica audible en el buzzer pasivo, y corroborar el estricto apego a las prohibiciones de alimentación (Req. P.17 y Q.18).
\end{tcolorbox}

\section{Materiales e Instrumentación a Solicitar en Bodega de Laboratorio}

Al presentarse en la ventanilla/bodega del laboratorio de la Escuela de Electrónica, solicite con su carné los siguientes equipos y accesorios:

\begin{enumerate}[leftmargin=*]
    \item \textbf{Fuente de Alimentación DC Regulada de Laboratorio (Doble Canal independiente):}
    \begin{itemize}
        \item \textbf{Canal 1 (CH1):} Ajustar a \textbf{$+5.00\text{ V}$} DC con límite de corriente fijado en $150\text{ mA}$ (alimenta el temporizador NE555, pines 8 y 4).
        \item \textbf{Canal 2 (CH2):} Ajustar a \textbf{$+8.50\text{ V}\text{ a }+9.00\text{ V}$} DC con límite de corriente en $100\text{ mA}$ (alimenta el pin 8 del LM358N para evitar saturación de riel en estados 12 a 15).
        \item \textbf{Tierra común:} Conectar los bornes negativos (negros) de ambos canales entre sí y referenciarlos al bus común de tierra (\texttt{GND}) de la protoboard.
    \end{itemize}
    \item \textbf{Multímetro Digital de Precisión (DMM):} Con un par de puntas tipo aguja/banana y al menos un par de terminales tipo caimán para mediciones fijas de tensión DC y resistencia ($\Omega$).
    \item \textbf{Osciloscopio Digital (2 Canales, $\ge 50\text{ MHz}$):} Con 2 sondas pasivas $10\times$ calibradas, puntas de agarre retráctil y cables de tierra con caimán corto.
    \item \textbf{Cables de Conexión de Laboratorio:}
    \begin{itemize}
        \item 3 cables Banana--Caimán (Rojo, Azul/Amarillo y Negro).
        \item 4 cables Caimán--Caimán para derivaciones de señal y puesta a tierra.
    \end{itemize}
    \item \textbf{Herramientas y Accesorios Propios:}
    \begin{itemize}
        \item Computadora portátil con cable USB para Arduino UNO, con el sketch \texttt{Logica.ino} listo y MATLAB abierto.
        \item Regla graduada o flexómetro (para certificar los $45\text{ cm}$ de la fibra óptica, Req. I.9).
        \item Destornillador plano/estrella pequeño (para ajustar la fijación de las torres M3 en la base láser).
    \end{itemize}
\end{enumerate}

\section{Fase 0: Inspección Previa y Conexión Segura (Pre-Flight Checks)}

\begin{itemize}[label=\squarebox]
    \item \textbf{Ensamble e Integridad Previa (Prueba de Banco):} El circuito físico de 7 etapas se encuentra 100\% ensamblado sobre la protoboard y la base láser. Tras la corrección de los lazos del LM358N, las pruebas de banco validaron que entrega los voltajes correctos de la escalera.
    \item \textbf{Verificación de Puentes del LM358N (Corrección de Voltaje):} Confirmar visualmente que los tres puentes del operacional dual estén firmemente insertados según el diagrama \texttt{puentes\_lm358.svg}: 1) Pin 1 con Pin 2 (buffer DAC Etapa 2), 2) Pin 6 con Pin 7 (buffer interfaz Etapa 3), y 3) Pin 1 con Pin 5 (acople directo inter-etapas).
    \item \textbf{Prohibición P.17 Verificada:} Confirmar visualmente que el pin \textbf{$+5\text{V}$ del Arduino UNO está totalmente desconectado} de la protoboard.
    \item \textbf{Tierra Compartida:} Unir firmemente el pin \texttt{GND} del Arduino UNO con el bus de tierra general (\texttt{GND}) de la protoboard y el negativo de la fuente de laboratorio.
    \item \textbf{Límites de Corriente:} Antes de activar la salida de la fuente, fijar el límite de corriente en $100\text{ mA}$ en ambos canales para proteger los circuitos integrados ante posibles cortocircuitos accidentales.
    \item \textbf{Tensiones de Encendido:} Activar primero CH1 ($+5.0\text{ V}$) y luego CH2 ($+9.0\text{ V}$). Observar en la fuente que el consumo global en reposo se mantenga por debajo de $40\text{ mA}$.
\end{itemize}

\section{Fase 1: Verificación de la Etapa 1 -- Microcontrolador Arduino (Req. C.3 y D.4)}

\begin{enumerate}[leftmargin=*]
    \item Abrir el Monitor Serial del Arduino IDE a \textbf{9600 baudios}.
    \item Enviar el comando \texttt{T1500} para configurar una permanencia cómoda de $1.5\text{ s}$ por nivel.
    \item \textbf{Comprobación de límites de retención (Req. C.3):}
    \begin{itemize}
        \item Probar \texttt{T1} ($1\text{ ms}$, límite inferior) $\to$ observar respuesta rápida en osciloscopio.
        \item Probar \texttt{T10000} ($10\text{ s}$, límite superior) $\to$ observar cambio pausado.
        \item Probar \texttt{T1500} para continuar con la toma ordenada de datos.
    \end{itemize}
    \item \textbf{Comprobación de patrones de MATLAB (Req. D.4):}
    \begin{itemize}
        \item Presionar \texttt{P}: Verificar que ejecuta el ciclo de la \textbf{señal periódica} (20 muestras).
        \item Presionar \texttt{S}: Verificar que ejecuta el ciclo de la \textbf{señal aperiódica} (18 muestras).
        \item Presionar \texttt{A}: Regresar al modo \textbf{escalera completa} (16 estados, 0 a 15).
    \end{itemize}
\end{enumerate}

\section{Fase 2: Verificación del DAC R--2R y Buffer LM358 (Req. E.5, F.6 y O.16)}

Colocar el sistema en modo manual enviando el código numérico correspondiente por el Monitor Serial (o utilizar el comando \texttt{M} para pausar). Con el multímetro digital referenciado a \texttt{GND}, medir la tensión continua en la salida de la Etapa 2 (Pin 1 del LM358N) y en la salida de la Interfaz Etapa 3 (Pin 7 del LM358N).

\begin{table}[h!]
\centering
\caption{Registro Experimental de Tensiones del DAC R--2R y Margen de Error (Req. O.16)}
\label{tab:dac_exp}
\footnotesize
\setlength{\tabcolsep}{4.5pt}
\renewcommand{\arraystretch}{1.12}
\begin{tabular}{ccccccc}
\toprule
\textbf{Dec} & \textbf{Binario} & $V_{\text{teo}}$ \textbf{(V)} & $V_{\text{sim}}$ \textbf{(V)} & $V_{\text{exp}}$ \textbf{Pin 1 (V)} & $\text{Error Relativo } E_{\text{rel}}$ \textbf{(\%)} & \textbf{Cumple O.16 ($< 10\%$)} \\
\midrule
0  & 0000 & 0.0000 & 0.0050 & \hspace{1.8cm} & Sat. Baja ($V_{OL} \approx 5\text{ mV}$) & [\quad] Sí \quad [\quad] No \\
1  & 0001 & 0.3125 & 0.3125 &                & & [\quad] Sí \quad [\quad] No \\
2  & 0010 & 0.6250 & 0.6250 &                & & [\quad] Sí \quad [\quad] No \\
3  & 0011 & 0.9375 & 0.9375 &                & & [\quad] Sí \quad [\quad] No \\
4  & 0100 & 1.2500 & 1.2500 &                & & [\quad] Sí \quad [\quad] No \\
5  & 0101 & 1.5625 & 1.5625 &                & & [\quad] Sí \quad [\quad] No \\
6  & 0110 & 1.8750 & 1.8750 &                & & [\quad] Sí \quad [\quad] No \\
7  & 0111 & 2.1875 & 2.1875 &                & & [\quad] Sí \quad [\quad] No \\
8  & 1000 & 2.5000 & 2.5000 &                & & [\quad] Sí \quad [\quad] No \\
9  & 1001 & 2.8125 & 2.8125 &                & & [\quad] Sí \quad [\quad] No \\
10 & 1010 & 3.1250 & 3.1250 &                & & [\quad] Sí \quad [\quad] No \\
11 & 1011 & 3.4375 & 3.4375 &                & & [\quad] Sí \quad [\quad] No \\
12 & 1100 & 3.7500 & 3.7500 &                & & [\quad] Sí \quad [\quad] No \\
13 & 1101 & 4.0625 & 4.0624 &                & & [\quad] Sí \quad [\quad] No \\
14 & 1110 & 4.3750 & 4.3749 &                & & [\quad] Sí \quad [\quad] No \\
15 & 1111 & 4.6875 & 4.6874 &                & & [\quad] Sí \quad [\quad] No \\
\bottomrule
\end{tabular}
\end{table}

\noindent\textbf{Fórmula de Cálculo del Error Relativo:}
\begin{equation*}
    E_{\text{rel}}(\%) = \left| \frac{V_{\text{exp}} - V_{\text{teo}}}{V_{\text{teo}}} \right| \times 100\%
\end{equation*}

\noindent\textbf{Prueba de Efecto de Carga (Req. F.6):}
\begin{itemize}[label=\squarebox]
    \item Fijar el código 15 ($4.6875\text{ V}$).
    \item Medir tensión en Pin 1 (salida DAC): $V_1 = \underline{\hspace{2cm}}\text{ V}$.
    \item Conectar y desconectar la carga de $150\ \Omega$ + LED en el Pin 7: Corroborar que $V_1$ en el Pin 1 se mantiene constante sin caerse (demuestra que la Etapa 3 aísla perfectamente la impedancia $R_{th} = 10\text{ k}\Omega$).
\end{itemize}

\section{Fase 3: Verificación del Generador y Enlace Óptico (Req. G.7, H.8 e I.9)}

\begin{itemize}[label=\squarebox]
    \item \textbf{Umbral de Conducción del LED Rojo:} Observar que para los códigos 0 a 5 ($V \le 1.56\text{ V}$), el LED permanece apagado. En el código 6 ($1.875\text{ V} > V_D \approx 1.80\text{ V}$), el LED comienza a emitir fotones visibles.
    \item \textbf{Corriente Máxima del LED (Código 15):} Medir la caída de tensión sobre $R_{driver} = 150\ \Omega$:
    \begin{equation*}
        V_{R_{drv}} = \underline{\hspace{2cm}}\text{ V} \implies I_{LED} = \frac{V_{R_{drv}}}{150\ \Omega} = \underline{\hspace{2cm}}\text{ mA} \quad (\text{Valor esperado: } 18.5\text{--}19.2\text{ mA})
    \end{equation*}
    \item \textbf{Longitud del Medio Óptico (Req. I.9):} Medir con la regla la fibra óptica plástica montada sobre las torres y ranuras:
    \begin{equation*}
        L_{\text{fibra}} = \underline{\hspace{2.5cm}}\text{ cm} \quad (\text{Debe cumplir estrictamente } 30\text{ cm} \le L \le 60\text{ cm})
    \end{equation*}
\end{itemize}

\section{Fase 4: Receptor Óptico y Oscilador Astable NE555 (Req. J.10, K.11 y L.12)}

Conectar el Canal 1 del osciloscopio en el Pin 3 del temporizador NE555 y el multímetro en medición de resistencia sobre los terminales de la LDR.

\begin{table}[h!]
\centering
\caption{Registro Experimental del Receptor LDR, Frecuencia y Audio Acústico}
\label{tab:555_exp}
\footnotesize
\setlength{\tabcolsep}{4.5pt}
\renewcommand{\arraystretch}{1.15}
\begin{tabular}{ccccccc}
\toprule
\textbf{Estado} & \textbf{Condición Óptica} & $R_{\text{LDR, sim}}$ & $R_{\text{LDR, exp}}$ \textbf{(k$\Omega$)} & $f_{\text{555, sim}}$ \textbf{(Hz)} & $f_{\text{555, exp}}$ \textbf{(Hz)} & \textbf{Tono Audible / Percepción} \\
\midrule
Código 0  & Oscuridad / LED Apagado & $40.0\text{ k}\Omega$ & & $176.7\text{ Hz}$ & & Tono Grave Profundo \\
Código 6  & Inicio Conducción LED   & $28.7\text{ k}\Omega$ & & $244.9\text{ Hz}$ & & Tono Medio-Bajo \\
Código 10 & Iluminación Intermedia  & $3.68\text{ k}\Omega$ & & $1711.6\text{ Hz}$ & & Tono Medio \\
Código 15 & Máxima Luz Conducida    & $1.58\text{ k}\Omega$ & & $3443.0\text{ Hz}$ & & Tono Agudo Intenso \\
\bottomrule
\end{tabular}
\end{table}

\noindent\textbf{Comprobaciones Críticas del Transductor y Acoplo AC:}
\begin{itemize}[label=\squarebox]
    \item \textbf{Prueba del Buzzer Pasivo (Req. L.12):} Desconectar momentáneamente el buzzer del circuito y conectarlo directamente a una fuente continua de $+5\text{ V}$ DC:
    \begin{itemize}
        \item Si emite un pitido continuo $\to$ \xmark{} \textbf{Es un buzzer activo (NO permitido)}.
        \item Si solo emite un "clic" mecánico seco al conectar y queda mudo $\to$ \cmark{} \textbf{Es 100\% pasivo (CUMPLE)}.
    \end{itemize}
    \item \textbf{Verificación del Acoplo AC ($C_{\text{acoplo}} = 10\ \mu\text{F}$):} Con el osciloscopio en acoplamiento DC conectado directamente sobre el buzzer, verificar que la onda cuadrada oscila simétricamente alrededor de $0\text{ V}$ ($V_{\text{offset}} \approx 0\text{ V}$), confirmando el filtrado del nivel continuo de reposo ($+2.5\text{ V}$).
\end{itemize}

\section{Checklist Rápido de Cumplimiento de Requerimientos para la Entrega}

\begin{table}[h!]
\centering
\footnotesize
\renewcommand{\arraystretch}{1.1}
\begin{tabular}{cp{10.5cm}c}
\toprule
\textbf{Req.} & \textbf{Descripción Resumida} & \textbf{Estado} \\
\midrule
\textbf{A.1 / B.2} & 7 etapas delimitadas en chasis físico y grabado láser. & [\quad] \\
\textbf{C.3 / D.4} & Patrones de 4 bits de longitud variable en MATLAB y Arduino ($1\text{ ms} \le T \le 10\text{ s}$). & [\quad] \\
\textbf{E.5}       & DAC con red R--2R ($10\text{ k}\Omega / 20\text{ k}\Omega$) + OpAmp activo 1A del LM358N. & [\quad] \\
\textbf{F.6}       & Etapa de interfaz de aislamiento exclusivamente con OpAmp 1B del LM358N. & [\quad] \\
\textbf{G.7 / H.8} & Generador óptico con resistencia $150\ \Omega$ y LED rojo justificado por Planck y PMMA. & [\quad] \\
\textbf{I.9 / J.10}& Fibra óptica plástica de $45\text{ cm}$ ($30\text{--}60\text{ cm}$) y fotorresistencia LDR pasiva. & [\quad] \\
\textbf{K.11 / L.12}& Oscilador comercial integrado NE555 astable y buzzer piezoeléctrico pasivo. & [\quad] \\
\textbf{M.13 / N.14}& Prototipado físico en LAIMI: Base en corte láser y kit 3D de 8 piezas atornilladas M3. & [\quad] \\
\textbf{\~N.15 / O.16}& Circuito montado en protoboard y error de conversión DAC menor al $10\%$. & [\quad] \\
\textbf{P.17 / Q.18}& Pin 5V de Arduino desconectado, sin baterías; alimentado por Fuente DC de Laboratorio. & [\quad] \\
\bottomrule
\end{tabular}
\end{table}

\vspace{0.3cm}
\noindent\rule{\textwidth}{0.4pt}
\begin{center}
\textbf{Firma de Integrantes:} \underline{\hspace{4.5cm}} \qquad \textbf{Visto Bueno Profesor/Asistente:} \underline{\hspace{4.5cm}}
\end{center}

\end{document}
